\documentclass[../../../include/open-logic-section]{subfiles}

\begin{document}

\olfileid[hi]{sth}{choice}{vitali}
\olsection{परिशिष्ट: विताली का विरोधाभास}

बानाख--टार्स्की निर्माण स्वीकार्य है या नहीं, इसका वास्तविक बोध पाने के
लिए हमें उसके \emph{प्रमाण} की जाँच करनी चाहिए। दुर्भाग्य से उसके लिए जितना
बीजगणित चाहिए, उतना हम यहाँ प्रस्तुत नहीं कर सकते। फिर भी निम्न परिणाम पर
ध्यान केंद्रित करके हम कुछ संक्षिप्त टिप्पणियाँ दे सकते हैं जो
बानाख--टार्स्की के प्रमाण पर कुछ प्रकाश डालेंगी।\footnote{अधिक पूर्ण निरूपण
के लिए \cite{Weston2003} या \cite{Wagon2016} देखें।}

\begin{thm}[विताली का विरोधाभास ($\ZFC$ में)]\ollabel{vitaliparadox}
किसी भी वृत्त को गणनीयतः अनेक टुकड़ों में विघटित किया जा सकता है, जिन्हें
(घूर्णन और स्थानांतरण द्वारा) फिर जोड़कर उस वृत्त की दो प्रतियाँ बनाई जा
सकती हैं।
\end{thm}

विताली के विरोधाभास को सिद्ध करना बानाख--टार्स्की विरोधाभास से कहीं सरल है।
हमने इसे ``विताली का विरोधाभास'' कहा है क्योंकि यह विताली के
\citeyear{Vitali1905} के अमापनीय समुच्चय के निर्माण से निष्पन्न होता है।
किंतु विताली और बानाख--टार्स्की विरोधाभासों के प्रमाणों के समुच्चय-सैद्धांतिक
पक्ष बहुत समान हैं। परिणामों में मुख्य अंतर केवल यह है कि बानाख--टार्स्की
\emph{परिमित} विघटन पर विचार करता है, जबकि विताली का विरोधाभास
\emph{गणनीयतः अनंत} विघटन पर। \citet{Weston2003} के शब्दों में, विताली का
विरोधाभास ``निश्चित ही बानाख--टार्स्की विरोधाभास जितना चकित करने वाला नहीं
है, किंतु वह दिखाता है कि ज्यामितीय विरोधाभास `सरल' स्थितियों में भी हो सकते
हैं।''

विताली का विरोधाभास एक द्विविम आकृति, वृत्त, से संबंधित है। इसलिए हम समतल
$\Real^2$ पर काम करेंगे। $\rotationsgroup$ उन (दक्षिणावर्त) घूर्णनों का
समुच्चय हो जो बिंदुओं को मूलबिंदु के चारों ओर $[0,2\pi)$ के बीच
\emph{परिमेय} रेडियन मानों से घुमाते हैं। $\rotationsgroup$ के बारे में कुछ
बीजगणितीय तथ्य ये हैं (यदि परिणाम का कथन समझ में न आए, तो प्रमाण उसका अर्थ
स्पष्ट कर देगा):

\begin{lem}\ollabel{rotationsgroupabelian}
$\rotationsgroup$ फलनों के संयोजन के अंतर्गत एक आबेली {समूह} बनाता है।
\end{lem}

\begin{proof}
$0$ रेडियन के घूर्णन को $0_{\rotationsgroup}$ लिखने पर यह
$\rotationsgroup$ का तत्समक अवयव है, क्योंकि किसी भी
$\rho\in\rotationsgroup$ के लिए
$\comp{0_{\rotationsgroup}}{\rho} = \comp{\rho}{0_{\rotationsgroup}} =
\rho$।

प्रत्येक अवयव का प्रतिलोम है। जहाँ $\rho\in\rotationsgroup$, $r$ रेडियन
से घुमाता है, वहाँ $\rho^{-1}\in\rotationsgroup$, $2\pi-r$ रेडियन से
घुमाता है, जिससे $\rho\circ\rho^{-1}=0_\rotationsgroup$।

संयोजन साहचर्यात्मक है: किन्हीं $\rho,\sigma,\tau\in\rotationsgroup$ के
लिए $\comp{\rho}{(\comp{\sigma}{\tau})}=
\comp{(\comp{\rho}{\sigma})}{\tau}$।

संयोजन क्रमविनिमेय है: किन्हीं $\rho,\sigma\in\rotationsgroup$ के लिए
$\comp{\rho}{\sigma}=\comp{\sigma}{\rho}$।
\end{proof}

वास्तव में हम अपने समूह $\rotationsgroup$ को दो भागों में बाँट सकते हैं और
फिर किसी भी आधे से पूरा समूह पुनः प्राप्त कर सकते हैं:

\begin{lem}\ollabel{disjointgroup}
$\rotationsgroup$ का दो असंयुक्त समुच्चयों $\rotationsgroup_1$ और
$\rotationsgroup_2$ में विभाजन है, जिनमें से दोनों $\rotationsgroup$ के
लिए आधार हैं।
\end{lem}

\begin{proof}
$\rotationsgroup_1$ में $[0,\pi)$ के परिमेय रेडियन मानों वाले घूर्णन हों;
$\rotationsgroup_2=\rotationsgroup\setminus\rotationsgroup_1$ रखें।
प्राथमिक बीजगणित से,
$\Setabs{\comp{\rho}{\rho}}{\rho \in \rotationsgroup_1} =
\rotationsgroup$। $\rotationsgroup_2$ के लिए समान परिणाम प्राप्त किया जा
सकता है।
\end{proof}

समूहों के बारे में इस तथ्य का उपयोग हम \olref{vitaliparadox} स्थापित करने
में करेंगे। $\onesphere$ इकाई वृत्त हो, अर्थात् समतल के मूलबिंदु से ठीक एक
इकाई दूर बिंदुओं का समुच्चय, अर्थात्
$\Setabs{\tuple{r,s}\in\Real^2}{\sqrt{r^2+s^2}=1}$। हम $\onesphere$ पर
निम्न संबंध पर विचार करके उसे भागों में बाँटेंगे:
\[
	r \sim s \emph{ तभी जब }(\exists \rho \in \rotationsgroup)\rho(r) = s.
\]
अर्थात् $\onesphere$ के बिंदु इस संबंध से ठीक तभी जुड़े हैं जब मूलबिंदु के
चारों ओर परिमेय मान वाले घूर्णन से एक से दूसरे तक पहुँचा जा सकता है। जैसा
अपेक्षित है:

\begin{lem}
$\sim$ एक तुल्यता संबंध है।
\end{lem}

\begin{proof}
\olref{rotationsgroupabelian} का उपयोग करके तुच्छ।
\end{proof}

अब हम चयन का उपयोग करके ऐसा समुच्चय $C$ प्राप्त करते हैं जिसमें $\sim$ के
अंतर्गत $\onesphere$ के प्रत्येक तुल्यता वर्ग से ठीक एक सदस्य हो। अर्थात्
हम तुल्यता वर्गों के समुच्चय पर किसी चयन फलन $f$ पर विचार करते हैं,
\footnote{चूँकि $\rotationsgroup$ !!{enumerable} है, $E$ का प्रत्येक
!!{element} !!{enumerable} है। चूँकि $\onesphere$ !!{nonenumerable} है,
\olref{vitalicover} और \olref[card-arithmetic][simp]{kappaunionkappasize}
से निष्पन्न होता है कि $E$ !!{nonenumerable} है। अतः यह \emph{अ}गणनीय
चयन का उपयोग है।}
\[
	E = \Setabs{\equivrep{r}{\sim}}{r \in \onesphere},
\]
और $C=\ran{f}$ रखें। प्रत्येक घूर्णन $\rho\in\rotationsgroup$ के लिए
समुच्चय $\funimage{\rho}{C}$ में $C$ के प्रत्येक बिंदु पर घूर्णन $\rho$
लगाने से प्राप्त बिंदु हैं। अगले दो परिणाम दिखाते हैं कि ये समुच्चय वृत्त को
पूर्णतः और बिना अतिव्यापन के ढकते हैं:

\begin{lem}\ollabel{vitalicover}
$\onesphere = \bigcup_{\rho \in \rotationsgroup} \funimage{\rho}{C}$.
\end{lem}

\begin{proof}
$s\in\onesphere$ को नियत करें; कोई $r\in C$ है ऐसा कि
$r\in\equivrep{s}{\sim}$, अर्थात् $r\sim s$, अर्थात् किसी
$\rho\in\rotationsgroup$ के लिए $\rho(r)=s$।
\end{proof}

\begin{lem}\ollabel{vitalinooverlap}
यदि $\rho_1\neq\rho_2$, तो
$\funimage{\rho_1}{C}\cap\funimage{\rho_2}{C}=\emptyset$।
\end{lem}

\begin{proof}
मान लें $s\in\funimage{\rho_1}{C}\cap\funimage{\rho_2}{C}$। अतः किन्हीं
$r_1,r_2\in C$ के लिए $s=\rho_1(r_1)=\rho_2(r_2)$। फलतः
$\rho^{-1}_2(\rho_1(r_1))=r_2$ और
$\comp{\rho_1}{\rho^{-1}_2}\in\rotationsgroup$, इसलिए $r_1\sim r_2$।
चूँकि $C$, $\sim$ के अंतर्गत प्रत्येक तुल्यता वर्ग से ठीक एक सदस्य चुनता
है, $r_1=r_2$। अतः $s=\rho_1(r_1)=\rho_2(r_1)$ और इसलिए
$\rho_1=\rho_2$।
\end{proof}

अब हम अपने पहले के बीजगणितीय तथ्यों को वृत्त पर लागू करते हैं:

\begin{lem}\ollabel{pseudobanachtarski}
$\onesphere$ का दो असंयुक्त समुच्चयों $D_1$ और $D_2$ में ऐसा विभाजन है कि
$D_1$ को गणनीयतः अनेक समुच्चयों में विभाजित किया जा सकता है, जिन्हें
घुमाकर $\onesphere$ की एक प्रति बनाई जा सकती है (और इसी प्रकार $D_2$ के
लिए)।
\end{lem}

\begin{proof}
\olref{disjointgroup} के $\rotationsgroup_1$ और $\rotationsgroup_2$ का
उपयोग करके रखें:
\begin{align*}
	D_{1} &= \bigcup_{\rho \in \rotationsgroup_1} \funimage{\rho}{C} & 
	D_{2} &= \bigcup_{\rho \in \rotationsgroup_2} \funimage{\rho}{C}
\end{align*}
\olref{vitalicover} से यह $\onesphere$ का विभाजन है, और
\olref{vitalinooverlap} से $D_1$ और $D_2$ असंयुक्त हैं। निर्माण से $D_1$
को गणनीयतः अनेक समुच्चयों, प्रत्येक $\rho\in R_1$ के लिए
$\funimage{\rho}{C}$, में विभाजित किया जा सकता है। और इन्हें घुमाकर
$\onesphere$ की प्रति बनाई जा सकती है, क्योंकि \olref{disjointgroup} और
\olref{vitalicover} से $\onesphere=\bigcup_{\rho\in
\rotationsgroup}\funimage{\rho}{C}=\bigcup_{\rho\in
\rotationsgroup_1}\funimage{(\comp{\rho}{\rho})}{C}$। यही तर्क $D_2$ पर
लागू होता है। \end{proof}\noindent इससे तत्काल विताली का विरोधाभास
निष्पन्न होता है। क्योंकि हम $\onesphere$ को गणनीयतः अनेक टुकड़ों (विभिन्न
$\funimage{\rho}{C}$) में बाँटकर और फिर उन्हें दृढ़ता से चलाकर ($D_1$ के
प्रत्येक टुकड़े को केवल घुमाकर, और $D_2$ के प्रत्येक टुकड़े को पहले
स्थानांतरित और फिर घुमाकर) $\onesphere$ से उसकी \emph{दो} प्रतियाँ बना सकते
हैं।

प्रमाण की रणनीति को दोहराएँ। हमने समतल पर घूर्णनों के समूह के बारे में कुछ
बीजगणितीय तथ्यों से आरंभ किया। इस समूह का उपयोग करके हमने $\onesphere$ को
तुल्यता वर्गों में विभाजित किया। फिर प्रत्येक वर्ग से अवयव चुनने के लिए चयन
का उपयोग करके हम एक ``विरोधाभास'' पर पहुँचे।

बानाख--टार्स्की सिद्ध करने के लिए हम ठीक यही रणनीति उपयोग करते हैं। मुख्य
अंतर यह है कि बानाख--टार्स्की सिद्ध करने में प्रयुक्त बीजगणितीय तथ्य विताली
के विरोधाभास में प्रयुक्त तथ्यों से पर्याप्त अधिक जटिल हैं। किंतु उन
बीजगणितीय तथ्यों का चयन से कोई संबंध नहीं है। हम उन्हें शीघ्रता से संक्षेपित
करेंगे।

बानाख--टार्स्की सिद्ध करने के लिए हम पहले \olref{disjointgroup} का एक
सदृश रूप स्थापित करते हैं: किसी भी \emph{मुक्त समूह} को चार टुकड़ों में
बाँटा जा सकता है, जिन्हें सहज रूप से ``इधर-उधर चलाकर'' हम पूरे समूह की दो
प्रतियाँ पुनः प्राप्त कर सकते हैं।\footnote{\emph{चार} टुकड़ों का उपयोग कर
सकने का तथ्य \cite{Robinson1947} का है। आधुनिक प्रमाण के लिए
\citet[Theorem 5.2]{Wagon2016} देखें। समूह के टुकड़ों को ``चलाने'' के वर्णन
में हम \citet[p.~3]{Weston2003} का अनुसरण करते हैं।} फिर हम दिखाते हैं कि
$\Real^3$ के मूलबिंदु के चारों ओर दो विशिष्ट घूर्णनों का उपयोग घूर्णनों का
एक मुक्त समूह $F$ उत्पन्न करने में किया जा सकता है।\footnote{
\citet[Theorem 2.1]{Wagon2016} देखें।} (अभी तक चयन नहीं।) अब हम गोले की सतह
के बिंदुओं को ``समान'' मानते हैं तभी जब $F$ के किसी घूर्णन से एक से दूसरा
प्राप्त किया जा सके। फिर हम ``समान'' बिंदुओं के प्रत्येक तुल्यता वर्ग से ठीक
एक बिंदु चुनने के लिए \emph{चयन का उपयोग करते हैं}।
\olref{pseudobanachtarski} की तरह $F$ के अपने विभाजन को गोले की सतह पर लागू
करके हम उस सतह को चार टुकड़ों में बाँटते हैं, जिन्हें ``इधर-उधर चलाकर'' गोले
की सतह की दो प्रतियाँ प्राप्त कर सकते हैं। इससे यह स्थापित होता है
\citep{Hausdorff1914}:

\begin{thm}[हाउसडॉर्फ़ का विरोधाभास ($\ZFC$ में)]
किसी भी गोले की सतह को परिमिततः अनेक टुकड़ों में विघटित किया जा सकता है,
जिन्हें (घूर्णन और स्थानांतरण द्वारा) फिर जोड़कर उस गोले की दो असंयुक्त
प्रतियाँ बनाई जा सकती हैं।
\end{thm}

पूर्ण बानाख--टार्स्की प्रमेय (जो केवल गोले की सतह ही नहीं, उसके आंतरिक भाग
से भी संबंधित है) प्राप्त करने के लिए कुछ और बीजगणितीय युक्तियाँ चाहिए।
किंतु सच कहें तो यह बीजगणितीय केक पर सजावट मात्र है। इसलिए वेस्टन लिखते हैं:
\begin{quote}	
  [\ldots] मुक्त समूहों का परिणाम बानाख--टार्स्की विरोधाभास के प्रमाण का
  \emph{मुख्य चरण} है। इस दृष्टि से बानाख--टार्स्की विरोधाभास $\Real^3$ के
  बारे में उतना कथन नहीं है, जितना वह [$\Real^3$ में स्थानांतरणों और
  घूर्णनों के] समूह की जटिलता के बारे में कथन है। \cite[p.~16]{Weston2003}
\end{quote}
अर्थात् हम \emph{परिमित} विघटन (बानाख--टार्स्की की तरह) दे सकते हैं या
\emph{गणनीयतः अनंत} विघटन (विताली के विरोधाभास की तरह), यह दो अथवा तीन
विमाओं में काम करने संबंधी कुछ समूह-सैद्धांतिक तथ्यों पर निर्भर करता है।

मानना होगा कि अंतिम प्रेक्षण \olref[banach]{sec} के अंत के चुटकुले को कुछ
बिगाड़ देता है। क्योंकि वह द्विविम है, यदि उन टुकड़ों को फिर व्यवस्थित करके
``Banach-Tarski Banach-Tarski'' बनाना हो, तो ``Banach-Tarski'' को गणनीयतः
\emph{अनंत} टुकड़ों में बाँटना होगा। चुटकुला सुधारने के लिए तीन विमाओं में
लिखना होगा। इसे हम पाठक के अभ्यास के लिए छोड़ते हैं।

एक अंतिम टिप्पणी। \olref[banach]{sec} में हमने कहा था कि प्राप्त होने वाले
गोले के ``टुकड़े'' \emph{मापनीय} नहीं हो सकते, बल्कि वे अचित्रणीय ``अनंत
बिखराव'' होने चाहिए। \olref{pseudobanachtarski} प्राप्त करने में चयन के
हमारे उपयोग के बारे में भी यही सत्य है। इसकी व्याख्या करना उपयोगी है।

फिर हमें कुछ पृष्ठभूमि की रूपरेखा देनी होगी (किंतु यह \emph{केवल} रूपरेखा
है; आप \emph{माप} पर किसी पाठ्यपुस्तक की प्रविष्टि देख सकते हैं)। समुच्चय
$X$ के लिए माप परिभाषित करना, $X$ पर किसी ``$\sigma$-बीजगणित'' के प्रत्येक
$E$ को मान $\mu(E)\in\Real$ देना है। यहाँ विवरण आवश्यक नहीं हैं, सिवाय इसके
कि फलन $\mu$ को गणनीय योज्यता के सिद्धांत का पालन करना चाहिए: असंयुक्त
समुच्चयों के गणनीय सम्मिलन का माप उनके अलग-अलग मापों का योग है, अर्थात् जब
$X_n$ असंयुक्त हों, तब
$\mu(\bigcup_{n<\omega}X_n)=\sum_{n<\omega}\mu(X_n)$। किसी समुच्चय को
``अमापनीय'' कहने का अर्थ है कि उसे कोई माप उपयुक्त रूप से नहीं दिया जा
सकता। अब पहले के अपने $\rotationsgroup$ का उपयोग करते हुए:

\begin{cor}[Vitali]
$\mu$ ऐसा माप हो कि $\mu(\onesphere)=1$, और यदि $X$ और $Y$ सर्वांगसम हों,
तो $\mu(X)=\mu(Y)$। तब सभी $\rho\in\rotationsgroup$ के लिए
$\funimage{\rho}{C}$ अमापनीय है।
\end{cor}

\begin{proof}
विरोधाभास द्वारा अन्यथा मान लें। अतः किन्हीं
$\sigma\in\rotationsgroup$ और $r\in\Real$ के लिए
$\mu(\funimage{\sigma}{C})=r$ मानें। किसी भी $\rho\in C$ के लिए
$\funimage{\rho}{C}$ और $\funimage{\sigma}{C}$ सर्वांगसम हैं, अतः किसी भी
$\rho\in C$ के लिए $\mu(\funimage{\rho}{C})=r$। \olref{vitalicover} और
\olref{vitalinooverlap} से
$\onesphere=\bigcup_{\rho\in\rotationsgroup}\funimage{\rho}{C}$ परस्पर
असंयुक्त समुच्चयों का गणनीय सम्मिलन है। अतः गणनीय योज्यता निर्देश देती है
कि $\mu(\onesphere)=1$ प्रत्येक $\funimage{\rho}{C}$ के मापों का योग है,
अर्थात्,
\[
	1 = \mu(\onesphere) = \sum_{\rho \in \rotationsgroup}\mu(\funimage{\rho}{C}) = \sum_{\rho \in \rotationsgroup}r
\]
किंतु यदि $r=0$, तो $\sum_{\rho\in\rotationsgroup}r=0$, और यदि $r>0$, तो
$\sum_{\rho\in\rotationsgroup}r=\infty$।
\end{proof}

\end{document}
